\documentclass[11pt,a4paper]{article}
\usepackage[utf8]{inputenc}\usepackage[T1]{fontenc}\usepackage[italian,english]{babel}
\usepackage{amsmath,amssymb,amsthm}\usepackage[margin=2.2cm]{geometry}\usepackage{longtable,booktabs,array,calc}
\usepackage{listings,xcolor}\usepackage{hyperref}\usepackage{url}
\providecommand{\tightlist}{\setlength{\itemsep}{0pt}\setlength{\parskip}{0pt}}
\lstset{basicstyle=\ttfamily\scriptsize,breaklines=true,frame=single,captionpos=b,columns=fullflexible,keepspaces=true,inputencoding=utf8,extendedchars=true}
\title{Proof Pursuit --- Problem 2, part 1\\ \large prove, decisioni degli agenti, codice verificato, letteratura}
\author{Team Proof Pursuit (Researcher: T. Tumini; Referee: gabundos; Creative: M. Renzi; gio3r3gio)}
\date{\today}
\begin{document}\maketitle\tableofcontents
\section*{Come leggere questo rapporto}
Per ogni cella: la richiesta ufficiale, lo stato dichiarato dal team (mai ``risolto'' senza revisione umana), la consegna
con la prova, poi la traccia delle decisioni degli agenti (Researcher: famiglia e motivo; Referee: verdetto ed errore
fatale; Creative: quando e perché è intervenuto) e il codice su cui il tentativo poggia con l'esito della riesecuzione
fidata. DIMOSTRATO, CITATO, VERIFICATO SU INTERVALLO FINITO e CONGETTURA restano distinti. L'architettura del sistema
è descritta nel README del repository.
\section{Problem 2 — Uphill paths on the hypercube}
\subsection{Cella 1 (1 punti) — stato: revisionato}
\subparagraph{\texorpdfstring{Parte 1 (C1) --- \(U(Q_3)\) and
\(U(Q_4)\)}{Parte 1 (C1) --- U(Q\_3) and U(Q\_4)}}\label{parte-1-c1-uq_3-and-uq_4}

\textbf{Punteggio:} 1 point · \textbf{Valutazione:} Checked instantly

The two smallest interesting cubes: \(Q_3\) has \(8\) vertices and
\(12\) edges, \(Q_4\) has \(16\) vertices and \(32\) edges. Determine
\(U(Q_3)\) and \(U(Q_4)\).

\textbf{Enunciato protetto dato agli agenti (state.json).} \texttt{Determine U(Q\_3) and U(Q\_4), each with an explicit optimal labelling (list of the 2\textasciicircum{}d vertices in increasing label order as 0/1 strings).}

\textbf{Evidenza registrata in STATUS.md.} \$U(Q\_3)=14\$, \$U(Q\_4)=34\$: prove a mano rilette passo per passo + verifiche esatte indipendenti; bozza in submission/parte-1.md

\subsubsection{Consegna (prova)}
\paragraph{Problema 2 --- Parte 1 (C1) --- bozza di
consegna}\label{problema-2-parte-1-c1-bozza-di-consegna}

\textbf{Stato dichiarato: RISOLTA (bozza revisionata).} Entrambe le
metà: \(U(Q_3)=14\) e \(U(Q_4)=34\), con valore, etichettatura esplicita
e prova completa a mano; calcoli esatti a conferma.

\subparagraph{1. Risultato e ambito}\label{risultato-e-ambito}

\(U(Q_3) = 14\) e \(U(Q_4) = 34\). Etichettatura ottima di \(Q_4\)
(ordine crescente di etichetta; due valli, \texttt{0000} e
\texttt{1111}):
\texttt{0000,\ 1111,\ 0111,\ 1101,\ 1110,\ 1011,\ 0001,\ 1000,\ 0010,\ 0100,\ 1001,\ 0110,\ 0011,\ 1010,\ 0101,\ 1100}.

Per \(Q_3\): Etichettatura ottima (ordine crescente di etichetta;
posizione \(i\) della stringa = coordinata \(i\); la convenzione è
irrilevante perché le permutazioni di coordinate sono automorfismi di
\(Q_3\)): \texttt{000,\ 100,\ 010,\ 110,\ 101,\ 011,\ 001,\ 111}.

\subparagraph{2. Dimostrazione}\label{dimostrazione}

\subparagraph{Result}\label{result}

\[U(Q_3) = 14.\]

An optimal labelling, as the list of the 8 vertices in increasing label
order (string position \(i\) = coordinate \(i\); the convention is
irrelevant for the count since \(\mathrm{Aut}(Q_3)\) contains all
coordinate permutations):

\[\texttt{000},\ \texttt{100},\ \texttt{010},\ \texttt{110},\ \texttt{101},\ \texttt{011},\ \texttt{001},\ \texttt{111}\]

i.e.~\(f(000)=1,\ f(100)=2,\ f(010)=3,\ f(110)=4,\ f(101)=5,\ f(011)=6,\ f(001)=7,\ f(111)=8\).

\subparagraph{Notation and the counting
identity}\label{notation-and-the-counting-identity}

For a labelling \(f\) of a graph \(G\), orient every edge from the
smaller to the larger label; this gives an acyclic orientation. Write
\(u \to v\) if \(uv\in E\) and \(f(u)<f(v)\), and
\(\deg^-(v)=\#\{u: u\to v\}\). Let \(p(v)\) be the number of uphill
paths whose last vertex is \(v\).

\textbf{Lemma 1 (recursion).}
\(p(v) = [\,v \text{ is a valley}\,] + \sum_{u\to v} p(u)\), and the
total number of uphill paths is \(\sum_v p(v)\).

\emph{Proof.} An uphill path ending at \(v\) has \(k=1\) (then it is
\((v)\) and \(v\) must be a valley; conversely a valley gives exactly
this one path) or \(k\ge2\); in the latter case deleting \(v\) gives an
uphill path ending at \(v_{k-1}\) with \(v_{k-1}\to v\), and conversely
appending \(v\) to any uphill path ending at some \(u\) with \(u\to v\)
gives an uphill path ending at \(v\) (the label condition \(f(u)<f(v)\)
is exactly \(u\to v\)). These correspondences are bijective, and the
sets for different \(u\) are disjoint (they differ in the penultimate
vertex). Every uphill path has a unique last vertex, hence the total is
\(\sum_v p(v)\). \(\square\)

\subparagraph{Upper bound: the labelling above has exactly 14 uphill
paths}\label{upper-bound-the-labelling-above-has-exactly-14-uphill-paths}

Neighbours in \(Q_3\): flip one coordinate. Compute \(p\) in increasing
label order using Lemma 1:

{\def\LTcaptype{none} % do not increment counter
\begin{longtable}[]{@{}llll@{}}
\toprule\noalign{}
vertex & label & lower neighbours & \(p\) \\
\midrule\noalign{}
\endhead
\bottomrule\noalign{}
\endlastfoot
000 & 1 & none (valley) & 1 \\
100 & 2 & 000 & 1 \\
010 & 3 & 000 & 1 \\
110 & 4 & 100, 010 & 1+1 = 2 \\
101 & 5 & 100 (001 has label 7, 111 has 8) & 1 \\
011 & 6 & 010 (001 has 7, 111 has 8) & 1 \\
001 & 7 & 000, 101, 011 & 1+1+1 = 3 \\
111 & 8 & 110, 101, 011 & 2+1+1 = 4 \\
\end{longtable}
}

Only 000 is a valley (every other vertex has a lower neighbour, as the
table shows). Total \(=1+1+1+2+1+1+3+4 = 14\). Hence \(U(Q_3)\le 14\).

\subparagraph{\texorpdfstring{Lower bound (hand proof): no labelling of
\(Q_3\) has \(\le 13\) uphill
paths}{Lower bound (hand proof): no labelling of Q\_3 has \textbackslash le 13 uphill paths}}\label{lower-bound-hand-proof-no-labelling-of-q_3-has-le-13-uphill-paths}

\textbf{Lemma 2.} For every labelling of any graph and every vertex
\(v\): \(p(v)\ge 1\), and \(p(v)\ge \deg^-(v)\).

\emph{Proof.} Starting from \(v\), repeatedly move to a neighbour with
smaller label while one exists; labels strictly decrease so this stops,
at a valley \(v_1\). Reversing the walk gives an uphill path ending at
\(v\), so \(p(v)\ge1\). Applying this to each \(u\) with \(u\to v\)
gives \(p(u)\ge1\), so by Lemma 1
\(p(v)\ge\sum_{u\to v}p(u)\ge\deg^-(v)\). \(\square\)

\textbf{Corollary 3.} Total
\(\ \ge \sum_v \max(1,\deg^-(v)) = \#\{\text{valleys}\} + \sum_{v \text{ not valley}}\deg^-(v) = \#\{\text{valleys}\} + |E|\),
since valleys are exactly the vertices with \(\deg^-=0\) and
\(\sum_v\deg^-(v)=|E|\) (each edge has exactly one head). For \(Q_3\),
\(|E|=12\) and the vertex with label 1 is always a valley, so the total
is \(\ge 13\).

\textbf{Proposition 4.} The total is never exactly 13 for \(Q_3\); hence
\(U(Q_3)\ge 14\).

\emph{Proof.} Suppose a labelling has total 13. Then every inequality in
Corollary 3 is an equality: (a) there is exactly one valley, and (b) for
every non-valley \(v\), \(p(v)=\deg^-(v)\),
i.e.~\(\sum_{u\to v}p(u)=\deg^-(v)\) with each \(p(u)\ge1\), so
\(p(u)=1\) for every \(u\to v\). By Lemma 2, \(p(u)=1\) forces
\(\deg^-(u)\le1\).

Let \(x\) be the vertex with label 8. All 3 neighbours of \(x\) are
lower, so \(\deg^-(x)=3\) and by (b) every \(a\in N(x)\) has
\(\deg^-(a)\le1\).

Let \(y\) be the vertex with label 7. At most one neighbour of \(y\)
(namely \(x\)) is higher, so \(\deg^-(y)\ge2\). If \(y\in N(x)\) we
would have \(\deg^-(y)\le1\), a contradiction; so \(y\notin N(x)\),
\(y\ne x\), hence \(y\) is the antipode \(\bar x\) of \(x\) (in \(Q_3\)
the only vertex at distance 3). Then \(N(y)\) is the set of the 3
vertices at distance 2 from \(x\), all lower than \(y\), so
\(\deg^-(y)=3\) and by (b) every \(b\in N(y)\) has \(\deg^-(b)\le1\).

\(N(x)\cup N(y)\) is the set \(S\) of the 6 vertices other than \(x,y\).
Count edges inside \(S\): \(Q_3\) has 12 edges, 3 are incident to \(x\),
3 to \(y\), and \(xy\) is not an edge, so exactly \(12-6=6\) edges have
both ends in \(S\). Each such edge is oriented towards its higher
endpoint, which lies in \(S\), so \(\sum_{s\in S}\deg^-(s)\ge 6\). But
each \(s\in S\) has \(\deg^-(s)\le1\), so
\(\sum_{s\in S}\deg^-(s)\le6\); therefore \(\deg^-(s)=1\) for every
\(s\in S\). Thus no vertex of \(S\) is a valley; \(x\) and \(y\) are not
valleys either (\(\deg^-\ge2\)). So the labelling has no valley,
contradicting the fact that the vertex with label 1 is a valley (all its
neighbours have larger labels). \(\square\)

Combining: \(U(Q_3)=14\), attained by the labelling above.

\subparagraph{Independent confirmation: exhaustive exact
computation}\label{independent-confirmation-exhaustive-exact-computation}

Script \texttt{q3\_esaustivo.py} (saved in \texttt{runs/p2\_q3/sandbox/}
and copied to \texttt{problema-2/certificati/}) enumerates \textbf{all}
\(8!=40320\) bijections \(V(Q_3)\to\{1,\dots,8\}\)
(\texttt{itertools.permutations}), with vertices encoded as integers
\(0..7\) (bit \(i\) = coordinate \(i\)) and adjacency = XOR with a
single bit. For each labelling it computes the number of uphill paths in
two independent ways --- (i) an explicit recursive DFS that enumerates
every uphill path from every valley, (ii) the DP of Lemma 1 processed in
increasing label order --- and \texttt{assert}s that the two agree. All
arithmetic is Python integers (exact; rigor = \texttt{exact}).

Output (Python 3.14.7, wall-clock 0.36 s):

\begin{verbatim}
etichettature esaminate: 40320
U(Q_3) = 14; etichettature ottime: 7104
distribuzione (cammini -> numero etichettature): [(14, 7104), (16, 14112), (17, 6720), (18, 4704), (19, 576), (20, 3648), (21, 1920), (22, 480), (23, 192), (24, 768), (26, 96)]
prima etichettatura ottima (ordine crescente di etichetta): ['000', '100', '010', '110', '101', '011', '001', '111']
\end{verbatim}

The distribution confirms Proposition 4 (no labelling with 13, and also
none with 15) and Corollary 3 (none below 13). Reproduce with
\texttt{.venv/bin/python\ problema-2/certificati/q3\_esaustivo.py}. A
second script \texttt{verifica\_etichettatura\_q3.py} checks only the
submitted labelling (bijection check, valleys, per-vertex \(p\), total
14) and prints the table above.

\subparagraph{\texorpdfstring{2b. Dimostrazione per
\(Q_4\)}{2b. Dimostrazione per Q\_4}}\label{b.-dimostrazione-per-q_4}

\paragraph{\texorpdfstring{\(U(Q_4)=34\)}{U(Q\_4)=34}}\label{uq_434}

\textbf{Conventions.} \(V(Q_4)=\{0,1\}^4\), \(u\sim v\) iff they differ
in exactly one coordinate; \(|V|=16\), \(|E|=32\), every vertex has
degree 4. In a 0/1 string the \(i\)-th character (from the left) is
coordinate \(i\); the convention is irrelevant for the count, since
every coordinate permutation is an automorphism of \(Q_4\) and
automorphisms preserve the number of uphill paths (they map valleys to
valleys and uphill paths to uphill paths bijectively when the labelling
is transported). For a labelling \(f\) orient each edge from the smaller
to the larger label (\(u\to v\) iff \(uv\in E\), \(f(u)<f(v)\));
\(\deg^-(v)=\#\{u:u\to v\}\), \(\deg^+(v)=\#\{w: v\to w\}\), so
\(\deg^-(v)+\deg^+(v)=4\). Let \(p(v)\) be the number of uphill paths
whose last vertex is \(v\).

\subparagraph{Preliminaries (verified claims, restated with proof for
self-containedness)}\label{preliminaries-verified-claims-restated-with-proof-for-self-containedness}

\textbf{Lemma 1.} \(p(v)=[v\text{ is a valley}]+\sum_{u\to v}p(u)\), and
the total number of uphill paths is \(\sum_v p(v)\).

\emph{Proof.} An uphill path ending at \(v\) either has length \(1\) ---
then it is \((v)\) and \(v\) is a valley, and conversely a valley \(v\)
gives exactly this path --- or has length \(\ge2\); deleting \(v\) gives
an uphill path ending at its penultimate vertex \(u\), with \(u\to v\),
and conversely appending \(v\) to an uphill path ending at \(u\) with
\(u\to v\) gives an uphill path ending at \(v\) (the increasing
condition is exactly \(f(u)<f(v)\)). These maps are mutually inverse,
and different \(u\) give disjoint sets (they differ in the penultimate
vertex). Each uphill path has a unique last vertex, hence the total is
\(\sum_v p(v)\). \(\square\)

\textbf{Lemma 2.} For every vertex \(v\): \(p(v)\ge1\) and
\(p(v)\ge\deg^-(v)\).

\emph{Proof.} From \(v\) repeatedly move to a neighbour with a smaller
label while one exists; the labels strictly decrease, so the walk stops,
at a vertex with no smaller neighbour, i.e.~a valley \(v_1\). The
reversed walk is an uphill path ending at \(v\), so \(p(v)\ge1\).
Applying this to every \(u\) with \(u\to v\) and using Lemma 1,
\(p(v)\ge\sum_{u\to v}p(u)\ge\sum_{u\to v}1=\deg^-(v)\). \(\square\)

\textbf{Corollary 3.} Total
\(\ \ge\sum_v\max(1,\deg^-(v))=\#\{\text{valleys}\}+|E|\).

\emph{Proof.} Termwise from Lemma 2 and Lemma 1. Valleys are exactly the
vertices with \(\deg^-=0\), and \(\sum_v\deg^-(v)=|E|\) because every
edge has exactly one head. So
\(\sum_v\max(1,\deg^-(v))=\#\{v:\deg^-(v)=0\}+\sum_{v:\deg^-(v)\ge1}\deg^-(v)=\#\{\text{valleys}\}+|E|\).
\(\square\)

The vertex with label \(1\) is always a valley (all its neighbours have
larger labels). Hence for \(Q_4\): total \(\ge 32+1=33\).

\subparagraph{\texorpdfstring{Lower bound: no labelling of \(Q_4\) has
exactly 33 uphill paths, hence
\(U(Q_4)\ge34\)}{Lower bound: no labelling of Q\_4 has exactly 33 uphill paths, hence U(Q\_4)\textbackslash ge34}}\label{lower-bound-no-labelling-of-q_4-has-exactly-33-uphill-paths-hence-uq_4ge34}

Suppose some labelling \(f\) of \(Q_4\) has total exactly \(33\).

\textbf{Step 1 (all inequalities are tight).} By Corollary 3,
\(33=\text{total}\ge\#\{\text{valleys}\}+32\ge 33\), so
\(\#\{\text{valleys}\}=1\) and \(\sum_v p(v)=\sum_v\max(1,\deg^-(v))\).
Since \(p(v)\ge\max(1,\deg^-(v))\) holds for every \(v\) (Lemma 2),
equality of the sums forces equality for every term:
\(p(v)=\max(1,\deg^-(v))\) for all \(v\). In particular, for every
non-valley \(v\) (i.e.~\(\deg^-(v)\ge1\)): \(p(v)=\deg^-(v)\).

\textbf{Step 2 (every in-neighbour of a non-valley has \(p=1\), hence
\(\deg^-\le1\)).} Let \(v\) be a non-valley. By Lemma 1 (with
\([v\text{ valley}]=0\)) and Step 1,
\[\sum_{u\to v}p(u)=p(v)=\deg^-(v)=\sum_{u\to v}1 .\] Each term
satisfies \(p(u)\ge1\) (Lemma 2), and there are \(\deg^-(v)\) terms on
each side, so \(p(u)=1\) for every \(u\to v\). By Lemma 2 again,
\(\deg^-(u)\le p(u)=1\).

\textbf{Step 3 (in-degrees lie in \(\{0,1,4\}\)).} Let \(u\) be any
vertex. If \(\deg^+(u)\ge1\), pick \(w\) with \(u\to w\); then \(w\) has
an in-neighbour, so \(w\) is not a valley, and Step 2 applied to \(v=w\)
gives \(\deg^-(u)\le1\). If \(\deg^+(u)=0\), then
\(\deg^-(u)=4-\deg^+(u)=4\). Hence \(\deg^-(u)\in\{0,1,4\}\) for every
vertex \(u\).

\textbf{Step 4 (counting edges).} Let \(s=\#\{u:\deg^-(u)=4\}\). The
vertices with \(\deg^-(u)=0\) are exactly the valleys, and there is
exactly one (Step 1). The remaining \(16-1-s=15-s\) vertices have
\(\deg^-(u)=1\). Summing in-degrees over all vertices, and using
\(\sum_u\deg^-(u)=|E|=32\):
\[4s+1\cdot(15-s)+0\cdot 1=32\quad\Longrightarrow\quad 3s=17,\] which
has no integer solution. Contradiction.

Therefore no labelling of \(Q_4\) has exactly \(33\) uphill paths;
combined with total \(\ge33\) (Corollary 3), every labelling has at
least \(34\) uphill paths: \(U(Q_4)\ge34\).

(\emph{Remark, not needed for the cell:} the same argument for \(Q_d\),
\(d\ge2\), gives total \(=|E|+1\) only if \((d-1)s=2^{d-1}(d-2)+1\) has
an integer solution; for \(d=3\) this is \(2s=5\), recovering
Proposition 4 of the \(Q_3\) proof.)

\subparagraph{Upper bound: an explicit labelling with exactly 34 uphill
paths}\label{upper-bound-an-explicit-labelling-with-exactly-34-uphill-paths}

Labelling \(f\), as the list of the 16 vertices in increasing label
order (label \(1\) first):

\[\texttt{0000},\ \texttt{1111},\ \texttt{0111},\ \texttt{1101},\ \texttt{1110},\ \texttt{1011},\ \texttt{0001},\ \texttt{1000},\ \texttt{0010},\ \texttt{0100},\ \texttt{1001},\ \texttt{0110},\ \texttt{0011},\ \texttt{1010},\ \texttt{0101},\ \texttt{1100}.\]

Structure: label \(1\) = weight \(0\); label \(2\) = weight \(4\);
labels \(3\)--\(6\) = the four vertices of weight \(3\); labels
\(7\)--\(10\) = the four vertices of weight \(1\); labels \(11\)--\(16\)
= the six vertices of weight \(2\). (This is a bijection: the sixteen
strings are pairwise distinct and are all of \(\{0,1\}^4\).)

Computation of \(p\) in increasing label order (Lemma 1). Neighbours of
a vertex of weight \(k\) have weights \(k\pm1\).

{\def\LTcaptype{none} % do not increment counter
\begin{longtable}[]{@{}
  >{\raggedright\arraybackslash}p{(\linewidth - 6\tabcolsep) * \real{0.2500}}
  >{\raggedright\arraybackslash}p{(\linewidth - 6\tabcolsep) * \real{0.2500}}
  >{\raggedright\arraybackslash}p{(\linewidth - 6\tabcolsep) * \real{0.2500}}
  >{\raggedright\arraybackslash}p{(\linewidth - 6\tabcolsep) * \real{0.2500}}@{}}
\toprule\noalign{}
\begin{minipage}[b]{\linewidth}\raggedright
vertex
\end{minipage} & \begin{minipage}[b]{\linewidth}\raggedright
label
\end{minipage} & \begin{minipage}[b]{\linewidth}\raggedright
lower neighbours
\end{minipage} & \begin{minipage}[b]{\linewidth}\raggedright
\(p\)
\end{minipage} \\
\midrule\noalign{}
\endhead
\bottomrule\noalign{}
\endlastfoot
0000 & 1 & none (its neighbours are the weight-1 vertices, labels 7--10)
--- \textbf{valley} & 1 \\
1111 & 2 & none (its neighbours are the weight-3 vertices, labels 3--6)
--- \textbf{valley} & 1 \\
0111 & 3 & 1111 (label 2); the other neighbours 0011, 0101, 0110 have
labels 13, 15, 12 & 1 \\
1101 & 4 & 1111; others 0101, 1001, 1100 have labels 15, 11, 16 & 1 \\
1110 & 5 & 1111; others 0110, 1010, 1100 have labels 12, 14, 16 & 1 \\
1011 & 6 & 1111; others 0011, 1001, 1010 have labels 13, 11, 14 & 1 \\
0001 & 7 & 0000 (label 1); others 1001, 0101, 0011 have labels 11, 15,
13 & 1 \\
1000 & 8 & 0000; others 1100, 1010, 1001 have labels 16, 14, 11 & 1 \\
0010 & 9 & 0000; others 1010, 0110, 0011 have labels 14, 12, 13 & 1 \\
0100 & 10 & 0000; others 1100, 0110, 0101 have labels 16, 12, 15 & 1 \\
1001 & 11 & 0001, 1000 (weight 1) and 1101, 1011 (weight 3): all four
lower, each with \(p=1\) & 4 \\
0110 & 12 & 0010, 0100, 1110, 0111 & 4 \\
0011 & 13 & 0001, 0010, 1011, 0111 & 4 \\
1010 & 14 & 1000, 0010, 1110, 1011 & 4 \\
0101 & 15 & 0001, 0100, 1101, 0111 & 4 \\
1100 & 16 & 1000, 0100, 1110, 1101 & 4 \\
\end{longtable}
}

Every weight-2 vertex has exactly two weight-1 and two weight-3
neighbours, all with labels \(\le10<11\), so all four neighbours are
lower and each has \(p=1\), giving \(p=4\). The valleys are exactly 0000
and 1111 (every other vertex has a lower neighbour, as the table shows).
Total \[\sum_v p(v)=1+1+4\cdot1+4\cdot1+6\cdot4=34 .\] Hence
\(U(Q_4)\le34\).

\subparagraph{Conclusion}\label{conclusion}

\[U(Q_4)=34,\] attained by the labelling above. (For comparison
\(U(Q_3)=14=|E|+2\) and \(U(Q_4)=34=|E|+2\).)

\subparagraph{Computational checks (exact integer arithmetic;
independent of the
proof)}\label{computational-checks-exact-integer-arithmetic-independent-of-the-proof}

\begin{enumerate}
\def\labelenumi{\arabic{enumi}.}
\item
  \texttt{verifica\_etichettatura\_q4.py} checks the submitted
  labelling: bijection, valleys, per-vertex \(p\) (the table above), and
  the total by two independent methods --- the DP of Lemma 1 in label
  order and an explicit recursive DFS enumerating every uphill path from
  every valley. Output: both totals \(=34\), valleys \(\{0000,1111\}\).
  Wall-clock 0.02 s (Python 3.14.7, macOS).
\item
  \texttt{q4\_branch\_and\_bound.py\ T} is an exact, symmetry-reduced
  search for labellings with \(<T\) uphill paths. Vertices are placed in
  increasing label order; when \(v\) is placed all its lower neighbours
  are already placed, so \(p(v)\) is fixed at that moment (Lemma 1).
  \emph{Symmetry reduction:} at each node only one candidate per orbit
  of the pointwise stabiliser (inside \(\mathrm{Aut}(Q_4)\), order
  \(384\), generated explicitly as coordinate permutation followed by
  XOR with a mask) of the already-placed set is tried; this is sound
  because if \(g\) fixes the placed vertices pointwise and \(g(v)=v'\),
  then \(f\mapsto f\circ g^{-1}\) maps labellings with prefix
  \((\dots,v)\) bijectively onto labellings with prefix \((\dots,v')\)
  and preserves the number of uphill paths. \emph{Lower bound at a node}
  (placed set \(S\), unplaced set \(U\),
  \(a(v)=\sum_{u\in N(v)\cap S}p(u)\), \(d(v)\) = number of
  in-neighbours of \(v\) inside \(U\) in the final orientation, so
  \(\sum_{v\in U}d(v)=E(U)\), the number of edges inside \(U\)): for
  \(v\in U\) all \(S\)-neighbours are lower and all unplaced
  in-neighbours have \(p\ge1\), so \(p(v)\ge\max(1,a(v)+d(v))\);
  minimising \(\sum_{v\in U}\max(1,a(v)+d(v))\) over nonnegative \(d\)
  with \(\sum d=E(U)\) gives
  \(\sum_{v\in U,a(v)\ge1}a(v)+\max(\#\{v\in U:a(v)=0\},E(U))\), so
  \(\text{LB}=\sum_{v\in S}p(v)+\sum_{v\in U,a(v)\ge1}a(v)+\max(\#\{a=0\},E(U))\)
  is a valid lower bound for every completion; a node is pruned iff
  \(\text{LB}\ge T\). Results: \(T=34\): \textbf{0 labellings found}, 47
  978 nodes, 0.3 s. \(T=35\): 19 525 440 labellings with exactly 34
  reached by the reduced search, 54 814 700 nodes, 156.8 s (shows the
  pruning is not over-aggressive at the optimum). Sanity check on
  \(Q_3\) (same code with \texttt{D\ =\ 3}): \(T=13\): 0 found;
  \(T=14\): 0 found (83 nodes); \(T=15\): 148 found, and
  \(148\cdot48=7104\) is exactly the known number of optimal labellings
  of \(Q_3\) (\(|\mathrm{Aut}(Q_3)|=48\)). This search is a confirmation
  only; the proof of \(U(Q_4)\ge34\) is the hand argument above.
\item
  \texttt{q4\_ricerca\_locale.py} (simulated annealing, 12 seeds, 200
  000 steps each, 12.8 s) found the value 34 in every run and never
  below; heuristic only, proves nothing, used to find the incumbent.
\end{enumerate}

\subparagraph{3. Verifica: istruzioni, dipendenze,
tempi}\label{verifica-istruzioni-dipendenze-tempi}

Tre programmi indipendenti, solo Python 3 standard (interi esatti),
ciascuno enumera tutte le \(8!=40320\) etichettature: -
\texttt{problema-2/certificati/q3\_esaustivo.py} --- DFS esplicita dei
cammini + programmazione dinamica di Lemma 1, con \texttt{assert} di
uguaglianza su ogni etichettatura. Tempo misurato: 0.37 s. -
\texttt{problema-2/certificati/q3\_verifica\_indipendente.py} ---
scritto separatamente, metodo diverso: costruzione esplicita livello per
livello di tutte le sequenze crescenti a partire dalle valli. Tempo
misurato: 0.20 s. -
\texttt{problema-2/certificati/verifica\_etichettatura\_q3.py} ---
controlla la sola etichettatura consegnata (biiezione, valli, \(p(v)\)
per vertice, totale 14). Comando:
\texttt{python3\ problema-2/certificati/q3\_verifica\_indipendente.py}
(e analoghi). Python 3.14.7, macOS. Tutti riportano: minimo 14, 7104
etichettature ottime, distribuzione senza valori 13 e 15.

Per \(Q_4\) (la prova del lower bound è a mano; il calcolo è solo
conferma): - \texttt{verifica\_etichettatura\_q4.py} (Researcher, due
metodi) e \texttt{q4\_verifica\_etichettatura\_indipendente.py} (scritto
a parte, costruzione esplicita dei cammini): entrambi 34 cammini, valli
\texttt{0000} e \texttt{1111}, biiezione verificata. -
\texttt{q4\_branch\_and\_bound.py\ SOGLIA}: branch-and-bound esatto con
riduzione per simmetria (384 automorfismi); con soglia 34 nessuna
etichettatura sotto 34 (47978 nodi, 0.3 s); con soglia 35 trova
etichettature con 34 (54.8M nodi, 157 s). La correttezza della riduzione
è argomentata nello script ma la cella non ne dipende. -
\texttt{q4\_ricerca\_locale.py}: simulated annealing, solo esplorazione
(12 seed, miglior valore 34 in ogni run).

\subparagraph{4. Fonti e contributo}\label{fonti-e-contributo}

Nessuna fonte esterna usata. L'argomento del lower bound generalizza
l'idea di IMO 2022 P6 (ogni vertice è fine di almeno un cammino;
\(p(v)\ge\deg^-(v)\)), aggiungendo l'ostruzione specifica di \(Q_3\)
(Proposizione 4). Prova prodotta dal Researcher automatico (Claude Fable
5.1, run \texttt{runs/p2\_q3/attempts/attempt\_001.json}) e revisionata
a mano passo per passo; verifica numerica indipendente scritta a parte.

\subparagraph{5. Limiti e osservazioni}\label{limiti-e-osservazioni}

\begin{itemize}
\tightlist
\item
  Nessuna parte irrisolta per C1.
\item
  \(U(Q_3)=|E|+2\) e \(U(Q_4)=|E|+2\). L'argomento di parità di \(Q_4\)
  (Step 3--4) si generalizza: il totale \(|E|+1\) è impossibile su
  \(Q_d\) quando \((d-1)\nmid 2^{d-1}(d-2)+1\), cioè per ogni \(d\) da 3
  a 10 almeno. Quindi \(U(Q_d)\ge|E|+2\) in quel range: bound debole
  rispetto ai 2368 noti per \(Q_9\), ma gratuito.
\item
  L'ottimo di \(Q_4\) ha DUE valli e realizza esattamente
  \(|E|+\#\text{valli}\): per i cubi grandi conviene cercare
  etichettature con poche valli in cui quasi tutti i vertici abbiano
  \(p(v)=\max(1,\deg^-(v))\).
\end{itemize}

\subsubsection{Decisioni degli agenti}
\paragraph{Run \texttt{p2\_c1}}
\paragraph{Run \texttt{p2\_q3}}
\begin{description}
\item[Iterazione 1 — attempt\_001]
\textbf{Researcher}: famiglia \texttt{exact\_computation}, sotto-obiettivo: Determine U(Q\_3) exactly (value + explicit optimal labelling in the required format), with certainty via exhaustive exact enumeration of all 8! labellings, plus an independent hand proof of the lower bound.. Stato dichiarato: \texttt{CELL\_SOLVED\_CANDIDATE}.
\emph{Perché questa scelta}: No prior attempts exist. The cell explicitly allows exhaustive exact computation with code and stated finite set; 8! = 40320 labellings is tiny (0.36 s). To make the result robust against a coding bug, two genuinely different counting routines are cross-checked on every labelling, and the lower bound is also proved by hand so the value does not rest on the code alone.
\textbf{Referee}: verdetto \texttt{(nessun verdetto)}.
\end{description}
\paragraph{Run \texttt{p2\_q4}}
\begin{description}
\item[Iterazione 1 — attempt\_001]
\textbf{Researcher}: famiglia \texttt{direct\_proof}, sotto-obiettivo: Determine U(Q\_4) exactly: prove the lower bound U(Q\_4) \textgreater{}= 34 by hand and exhibit an explicit labelling with exactly 34 uphill paths (route (b) of the cell), with an independent exact computer confirmation.. Stato dichiarato: \texttt{CELL\_SOLVED\_CANDIDATE}.
\emph{Perché questa scelta}: No failed attempts recorded. Simulated annealing (12 seeds) never went below 34 = |E| + 2, while the verified bound gives 33; the gap of 1 suggested a tightness analysis of the verified bound |E| + \#valleys, exactly as Proposition 4 did for Q\_3. The tightness analysis turned out to close with a pure counting argument (no case analysis), so I chose route (b) (hand proof + explicit labelling), and added a symmetry-reduced exact search only as an independent check, not as the proof.
\textbf{Referee}: verdetto \texttt{(nessun verdetto)}.
\end{description}

\subsubsection{Codice su cui poggia il tentativo}
Nessun codice nei tentativi.
\begin{lstlisting}[caption={attempt\_001 code\_1 [python, rigore exact] — Exhaustive exact enumeration of all 8! = 40320 labellings of Q\_3; two independent counts ( — verdetto del Referee: (nessun verdetto)}]
"""
Calcolo esatto di U(Q_3) per enumerazione esaustiva di tutte le 8! = 40320 etichettature.

Due conteggi indipendenti per ogni etichettatura:
  1. DFS: enumera esplicitamente ogni cammino in salita (v_1,...,v_k) partendo da ogni valle.
  2. DP:  p(v) = [v valle] + somma_{u vicino, f(u)<f(v)} p(u);  totale = somma_v p(v).
Aritmetica: solo interi Python (esatta). Insieme coperto: TUTTE le biiezioni V(Q_3) -> {1..8}.
"""
from itertools import permutations
from collections import Counter
import time

D = 3
VERTICI = list(range(2 ** D))                      # vertice = intero, bit i = coordinata i
VICINI = [[v ^ (1 << i) for i in range(D)] for v in VERTICI]


def stringa(v):
    """Vertice intero -> stringa 0/1 di lunghezza D (bit 0 a sinistra)."""
    return "".join(str((v >> i) & 1) for i in range(D))


def valli(etichetta):
    """Vertici con tutti i vicini di etichetta maggiore."""
    return [v for v in VERTICI if all(etichetta[w] > etichetta[v] for w in VICINI[v])]


def conta_dfs(etichetta):
    """Conteggio 1: enumera esplicitamente ogni cammino in salita."""
    def cammini_da(v):
        # il cammino che termina in v conta 1, piu' tutte le estensioni ai vicini maggiori
        return 1 + sum(cammini_da(w) for w in VICINI[v] if etichetta[w] > etichetta[v])
    return sum(cammini_da(v) for v in valli(etichetta))


def conta_dp(etichetta):
    """Conteggio 2: programmazione dinamica in ordine crescente di etichetta."""
    ordine = sorted(VERTICI, key=lambda v: etichetta[v])
    p = {}
    for v in ordine:
        minori = [w for w in VICINI[v] if etichetta[w] < etichetta[v]]
        p[v] = (1 if not minori else 0) + sum(p[w] for w in minori)
    return sum(p.values())


def esplora_tutte():
    """Enumera tutte le 8! etichettature; restituisce (minimo, distribuzione, ottime)."""
    distribuzione = Counter()
    ottime = []
    minimo = None
    for perm in permutations(range(1, 2 ** D + 1)):
        etichetta = list(perm)                     # etichetta[v] = f(v)
        n_dfs = conta_dfs(etichetta)
        n_dp = conta_dp(etichetta)
        assert n_dfs == n_dp, (etichetta, n_dfs, n_dp)
        distribuzione[n_dfs] += 1
        if minimo is None or n_dfs < minimo:
            minimo, ottime = n_dfs, [etichetta]
        elif n_dfs == minimo:
            ottime.append(etichetta)
    return minimo, distribuzione, ottime


def lista_ordinata(etichetta):
    """Formato di consegna: vertici in ordine crescente di etichetta, come stringhe 0/1."""
    return [stringa(v) for v in sorted(VERTICI, key=lambda v: etichetta[v])]


if __name__ == "__main__":
    inizio = time.time()
    minimo, distribuzione, ottime = esplora_tutte()
    print(f"etichettature esaminate: {sum(distribuzione.values())}")
    print(f"U(Q_3) = {minimo}; etichettature ottime: {len(ottime)}")
    print("distribuzione (cammini -> numero etichettature):", sorted(distribuzione.items()))
    print("prima etichettatura ottima (ordine crescente di etichetta):", lista_ordinata(ottime[0]))
    print(f"tempo: {time.time() - inizio:.2f} s")

\end{lstlisting}
\begin{lstlisting}[caption={attempt\_001 code\_2 [python, rigore exact] — Independent check of the single submitted labelling (bijection check, valleys, per-vertex  — verdetto del Referee: (nessun verdetto)}]
"""
Verifica indipendente della singola etichettatura consegnata per Q_3:
stampa valli, p(v) per ogni vertice e il totale dei cammini in salita (interi esatti).
"""
D = 3
CONSEGNA = ['000', '100', '010', '110', '101', '011', '001', '111']   # ordine crescente di etichetta


def vicini(s):
    """Stringhe 0/1 che differiscono da s in esattamente una coordinata."""
    return [s[:i] + ('1' if s[i] == '0' else '0') + s[i + 1:] for i in range(D)]


def verifica(consegna):
    """Calcola p(v) in ordine di etichetta e restituisce (valli, p, totale)."""
    assert sorted(consegna) == sorted(format(k, f'0{D}b') for k in range(2 ** D)), "non e' una biiezione"
    etichetta = {s: k + 1 for k, s in enumerate(consegna)}
    p, valli = {}, []
    for s in consegna:
        minori = [w for w in vicini(s) if etichetta[w] < etichetta[s]]
        if not minori:
            valli.append(s)
        p[s] = (1 if not minori else 0) + sum(p[w] for w in minori)
    return valli, p, sum(p.values())


if __name__ == "__main__":
    valli, p, totale = verifica(CONSEGNA)
    print("valli:", valli)
    for s in CONSEGNA:
        print(f"  {s}  etichetta {CONSEGNA.index(s)+1}  p = {p[s]}")
    print("cammini in salita totali:", totale)

\end{lstlisting}
\begin{lstlisting}[caption={attempt\_001 code\_1 [python, rigore exact] — Exact verification of the submitted labelling of Q\_4: bijection check, valleys, per-vertex — verdetto del Referee: (nessun verdetto)}]
"""Verifica ESATTA (interi Python) dell'etichettatura candidata di Q_4 con 34 cammini in salita.
Due metodi indipendenti: (i) DP di Lemma 1 in ordine di etichetta; (ii) enumerazione esplicita
(DFS) di tutti i cammini in salita da ogni valle. Stampa la tabella p(v) per la consegna."""
D = 4
N = 1 << D
# Convenzione: carattere i-esimo della stringa (da sinistra) = coordinata i; bit i dell'intero.
ETICHETTATURA = ['0000', '1111', '0111', '1101', '1110', '1011', '0001', '1000',
                 '0010', '0100', '1001', '0110', '0011', '1010', '0101', '1100']


def da_stringa(s):
    """Converte la stringa 0/1 (carattere i = coordinata i) nell'intero con bit i = coordinata i."""
    return sum(int(c) << i for i, c in enumerate(s))


def vicini(v):
    """I d vicini di v in Q_d: si cambia esattamente una coordinata."""
    return [v ^ (1 << i) for i in range(D)]


def conta_con_dp(ordine, etichetta):
    """Metodo (i): p(v) = [v valle] + somma p(u) sui vicini u con etichetta minore."""
    p = {}
    righe = []
    for v in ordine:
        minori = [u for u in vicini(v) if etichetta[u] < etichetta[v]]
        p[v] = 1 if not minori else sum(p[u] for u in minori)
        righe.append((format(v, f'0{D}b')[::-1], etichetta[v], [format(u, f'0{D}b')[::-1] for u in minori], p[v]))
    return sum(p.values()), righe


def conta_con_dfs(etichetta):
    """Metodo (ii): enumera esplicitamente ogni cammino in salita partendo da ogni valle."""
    def estendi(v):
        totale = 1  # il cammino che termina qui
        for w in vicini(v):
            if etichetta[w] > etichetta[v]:
                totale += estendi(w)
        return totale
    valli = [v for v in range(N) if all(etichetta[w] > etichetta[v] for w in vicini(v))]
    return sum(estendi(v) for v in valli), valli


def main():
    ordine = [da_stringa(s) for s in ETICHETTATURA]
    assert sorted(ordine) == list(range(N)), "non e' una biiezione"
    etichetta = {v: i + 1 for i, v in enumerate(ordine)}
    totale_dp, righe = conta_con_dp(ordine, etichetta)
    totale_dfs, valli = conta_con_dfs(etichetta)
    for stringa, lab, minori, p in righe:
        print(f"{stringa}  etichetta {lab:2d}  vicini minori {minori}  p = {p}")
    print("valli:", [format(v, f'0{D}b')[::-1] for v in valli])
    print("totale (DP di Lemma 1):", totale_dp)
    print("totale (DFS esplicita):", totale_dfs)
    assert totale_dp == totale_dfs == 34


if __name__ == "__main__":
    main()

\end{lstlisting}
\begin{lstlisting}[caption={attempt\_001 code\_2 [python, rigore exact] — Independent exact confirmation of the lower bound: symmetry-reduced branch-and-bound over  — verdetto del Referee: (nessun verdetto)}]
"""Branch-and-bound ESATTO (interi) su Q_4: cerca un'etichettatura con < SOGLIA cammini in salita.
Si piazzano i vertici in ordine crescente di etichetta; p(v) e' determinato al piazzamento (Lemma 1).
Riduzione per simmetria: a ogni nodo si prova un solo candidato per orbita dello stabilizzatore
puntuale (in Aut(Q_4), ordine 384) dell'insieme gia' piazzato. Bound inferiore dimostrato nel testo.
Conferma indipendente della prova a mano; NON e' la prova principale."""
import sys, time
from itertools import permutations

D = 4
N = 1 << D
VICINI = [[v ^ (1 << i) for i in range(D)] for v in range(N)]


def automorfismi():
    """Tutti i 384 automorfismi di Q_4 come tuple immagine: v -> perm(coordinate) XOR maschera."""
    gruppo = []
    for perm in permutations(range(D)):
        for maschera in range(N):
            immagine = tuple(sum(((v >> i) & 1) << perm[i] for i in range(D)) ^ maschera for v in range(N))
            gruppo.append(immagine)
    return gruppo


GRUPPO = automorfismi()


def bound_inferiore(p, piazzati, non_piazzati):
    """Bound: per v non piazzato p(v) >= max(1, a(v) + d(v)), a(v) = somma p sui vicini piazzati,
    d(v) = archi entranti da vertici non piazzati (somma dei d(v) = archi interni E(U)).
    Minimizzando sui d(v) a somma fissata: somma_{a>=1} a(v) + max(|U0|, E(U)), U0 = {a(v)=0}."""
    totale = sum(p[v] for v in piazzati)
    zeri, archi_interni = 0, 0
    for v in non_piazzati:
        a = sum(p[u] for u in VICINI[v] if u in piazzati)
        if a == 0:
            zeri += 1
        else:
            totale += a
        archi_interni += sum(1 for u in VICINI[v] if u not in piazzati)
    return totale + max(zeri, archi_interni // 2)


def rappresentanti(non_piazzati, stabilizzatore):
    """Un candidato per orbita dello stabilizzatore sui vertici non piazzati."""
    visti, scelti = set(), []
    for v in sorted(non_piazzati):
        if v in visti:
            continue
        scelti.append(v)
        visti.update(g[v] for g in stabilizzatore)
    return scelti


class Ricerca:
    def __init__(self, soglia):
        self.soglia = soglia
        self.nodi = 0
        self.trovate = []

    def esplora(self, ordine, p, piazzati, stabilizzatore):
        """Espande un nodo: prova ogni rappresentante d'orbita come prossimo vertice."""
        self.nodi += 1
        non_piazzati = [v for v in range(N) if v not in piazzati]
        if not non_piazzati:
            self.trovate.append(list(ordine))
            return
        if bound_inferiore(p, piazzati, non_piazzati) >= self.soglia:
            return
        for v in rappresentanti(non_piazzati, stabilizzatore):
            minori = [u for u in VICINI[v] if u in piazzati]
            p[v] = 1 if not minori else sum(p[u] for u in minori)
            piazzati.add(v)
            ordine.append(v)
            self.esplora(ordine, p, piazzati, [g for g in stabilizzatore if g[v] == v])
            ordine.pop()
            piazzati.discard(v)
            del p[v]


def main():
    soglia = int(sys.argv[1]) if len(sys.argv) > 1 else 34
    inizio = time.time()
    ricerca = Ricerca(soglia)
    ricerca.esplora([], {}, set(), GRUPPO)
    print(f"soglia {soglia}: etichettature con < {soglia} cammini trovate: {len(ricerca.trovate)}")
    for ordine in ricerca.trovate[:5]:
        print("  ", [format(v, f'0{D}b')[::-1] for v in ordine])
    print(f"nodi visitati: {ricerca.nodi}; tempo: {time.time() - inizio:.1f} s")


if __name__ == "__main__":
    main()

\end{lstlisting}
\begin{lstlisting}[caption={attempt\_001 code\_3 [python, rigore float\_exploration\_only] — Heuristic simulated annealing to find an incumbent labelling (upper bound candidate). 12 s — verdetto del Referee: (nessun verdetto)}]
"""Ricerca locale (simulated annealing) di etichettature di Q_4 con pochi cammini in salita.
Serve solo a trovare un buon candidato (upper bound): NON dimostra nulla.
Aritmetica: interi esatti (il conteggio), float solo nella regola di accettazione di SA."""
import random, sys, math

D = 4
N = 1 << D
VICINI = [[v ^ (1 << i) for i in range(D)] for v in range(N)]


def conta_cammini(ordine):
    """Numero di cammini in salita dell'etichettatura data come lista di vertici
    in ordine crescente di etichetta (Lemma 1: p(v) = [valle] + somma p(u) sui vicini minori)."""
    etichetta = [0] * N
    for pos, v in enumerate(ordine):
        etichetta[v] = pos
    p = [0] * N
    totale = 0
    for v in ordine:
        minori = [u for u in VICINI[v] if etichetta[u] < etichetta[v]]
        p[v] = 1 if not minori else sum(p[u] for u in minori)
        totale += p[v]
    return totale


def ricottura(seme, passi=200000):
    """Una corsa di simulated annealing con mosse di scambio e di spostamento."""
    rng = random.Random(seme)
    ordine = list(range(N))
    rng.shuffle(ordine)
    valore = conta_cammini(ordine)
    migliore, migliore_ordine = valore, ordine[:]
    for passo in range(passi):
        temperatura = 2.0 * (1 - passo / passi) + 0.05
        nuovo = ordine[:]
        i, j = rng.randrange(N), rng.randrange(N)
        if rng.random() < 0.5:
            nuovo[i], nuovo[j] = nuovo[j], nuovo[i]
        else:
            nuovo.insert(j, nuovo.pop(i))
        nuovo_valore = conta_cammini(nuovo)
        if nuovo_valore <= valore or rng.random() < math.exp((valore - nuovo_valore) / temperatura):
            ordine, valore = nuovo, nuovo_valore
            if valore < migliore:
                migliore, migliore_ordine = valore, ordine[:]
    return migliore, migliore_ordine


def main():
    corse = int(sys.argv[1]) if len(sys.argv) > 1 else 20
    globale, globale_ordine = None, None
    for seme in range(corse):
        valore, ordine = ricottura(seme)
        if globale is None or valore < globale:
            globale, globale_ordine = valore, ordine
        print(f"seme {seme}: {valore} (migliore finora {globale})", flush=True)
    stringhe = [format(v, f"0{D}b") for v in globale_ordine]
    print("MIGLIORE:", globale, stringhe)


if __name__ == "__main__":
    main()

\end{lstlisting}

\section{arXiv literature consulted}
\begin{thebibliography}{99}
\bibitem{1412.3893v1} Ian E. Ochs, Michael M. Desai. \emph{The competition between simple and complex evolutionary trajectories in asexual populations}. arXiv:1412.3893v1 (2014). \url{http://arxiv.org/abs/1412.3893v1}
\end{thebibliography}
\end{document}
